\chapter{Monte Carlo}\label{ch:qm:monte-carlo}

The overnight batch values 40\,000 path-dependent trades with 100\,000 paths each. The \omterm{def:qm:estimation:estimator}{standard error} on the book is 80\,000 dollars, and the risk manager wants it under 20\,000 by tomorrow.
Monte Carlo error falls like one over the square root of the number of paths, so a factor of four in error is a factor of sixteen in paths, and sixteen times the machines is not on offer. The
alternative is to make each path worth more. This chapter builds the \omterm{def:qm:estimation:estimator}{estimator} and its generators, then the variance-reduction techniques that change the constant (\omterm{def:qm:monte-carlo:vr}{antithetic variates},
\omterm{def:qm:monte-carlo:vr}{control variates}, stratification, \omterm{def:qm:monte-carlo:is}{importance sampling}), the low-discrepancy points that change the rate, and the discretisation and multilevel methods that decide what a path costs.

\section{Estimators and generators}

\begin{definition}[Monte Carlo method, variance reduction]\label{def:qm:monte-carlo:mc}
The \emph{Monte Carlo method}\index{Monte Carlo method} estimates $\mu = \E[f(X)]$ by the average $\hat\mu_n = \frac1n\sum_{i=1}^nf(X_i)$ of independent draws. \emph{Variance
reduction}\index{variance reduction} is any change of \omterm{def:qm:estimation:estimator}{estimator} that keeps $\E[\hat\mu_n] = \mu$ (or makes the bias negligible) and lowers its variance at equal cost.
\end{definition}

\begin{proposition}[Error and cost of the Monte Carlo estimator]\label{prop:qm:monte-carlo:clt}
If $\sigma^2 = \Var(f(X)) < \infty$, then $\sqrt n(\hat\mu_n - \mu) \Rightarrow \mathcal N(0, \sigma^2)$, so $\hat\mu_n \pm 1.96\,s_n/\sqrt n$ is an asymptotic 95\% interval. If one draw costs $c$, the
cost of reaching \omterm{def:qm:estimation:estimator}{standard error} $\epsilon$ is $c\sigma^2/\epsilon^2$: of two \omterm{def:qm:estimation:estimator}{unbiased estimators}, the better is the one with the smaller product of variance and cost per draw.
\end{proposition}

\begin{proof}
The first statement is the central limit theorem for independent identically distributed draws, with Slutsky's lemma for the estimated $s_n$. The second: $\sigma/\sqrt n = \epsilon$ needs $n =
\sigma^2/\epsilon^2$ draws.
\end{proof}

The rate $n^{-1/2}$ does not depend on the dimension, which is why Monte Carlo prices 64-date averages and 30-asset baskets that no grid can reach; the constant $\sigma^2c$ is where the work
is. On the batch's representative trade, an arithmetic-average call on 64 fixings of a \omterm{def:qm:stochastic-differential-equations:gbm}{geometric Brownian motion} ($S_0 = K = 100$, $r = 3\%$, $\sigma = 25\%$, one year), $2^{14}$ paths give
6.495 with a \omterm{def:qm:estimation:estimator}{standard error} of 0.077; the price is 6.4788.

\begin{definition}[Pseudo-random and counter-based generators, inverse transform sampling]\label{def:qm:monte-carlo:prng}
A \emph{pseudo-random number generator}\index{pseudo-random number generator} is a deterministic recurrence $s_{k+1} = F(s_k)$, $u_k = G(s_k)$ whose outputs pass statistical tests of
independent uniformity; the Mersenne Twister (Matsumoto and Nishimura, 1998) and PCG (O'Neill, 2014; NumPy's default) are examples. A \emph{counter-based generator}\index{counter-based
generator} computes $u = G_k(c)$ directly from a counter $c$ and a key $k$ with a bijection $G_k$, so that any draw can be produced without its predecessors; Philox (Salmon, Moraes, Dror
and Shaw, 2011) is one. \emph{Inverse transform sampling}\index{inverse transform sampling} turns a uniform $U$ into a draw $F^{-1}(U)$ of the distribution function $F$.
\end{definition}

A \omterm{def:qm:monte-carlo:prng}{counter-based generator} is what makes a distributed batch reproducible. The engine gives path $j$ of trade $i$ the counter $(b, j, i, 0)$ for its $b$-th block of four numbers, so
the path's numbers do not depend on how the paths are split among machines, on the order in which they run, or on how many there are: rerunning one path to debug a trade reproduces it bit for
bit. The kit's Philox4x64-10 reproduces NumPy's implementation block for block, in Python, C++20 and Rust (the C++ version is listed in \cref{tut:qm:monte-carlo:asian}). Normals come from the inverse transform with
Acklam's rational approximation of $\Phi^{-1}$ (relative error below $1.15 \times 10^{-9}$), which, unlike Box--Muller, maps one uniform to one normal: the property low-discrepancy points need.

\section{Variance reduction and importance sampling}

\begin{definition}[Antithetic variates, control variate, stratified sampling]\label{def:qm:monte-carlo:vr}
\emph{Antithetic variates}\index{antithetic variates} pair each draw $Z$ with $-Z$ (each uniform $U$ with $1 - U$) and average $\frac12(f(Z) + f(-Z))$. A \emph{control variate}\index{control
variate} is a variable $C$ with known mean, used in $\hat\mu_\beta = \frac1n\sum_i(Y_i - \beta(C_i - \E C))$. \emph{Stratified sampling}\index{stratified sampling} partitions the space of an input into
strata of known probability $p_s$ and draws a fixed number $n_s = np_s$ in each.
\end{definition}

\begin{proposition}[Optimal control-variate coefficient]\label{prop:qm:monte-carlo:cv}
$\Var(Y - \beta(C - \E C))$ is minimised at $\beta^* = \Cov(Y, C)/\Var(C)$, where it equals $(1 - \rho^2)\Var(Y)$ with $\rho = \mathrm{Corr}(Y, C)$.
\end{proposition}

\begin{proof}
$\Var(Y - \beta C) = \Var Y - 2\beta\Cov(Y, C) + \beta^2\Var C$ is a quadratic in $\beta$ minimised at $\beta^*$; substituting gives $\Var Y - \Cov(Y, C)^2/\Var C = (1 - \rho^2)\Var Y$.
\end{proof}

For the average-price call, the geometric average $G = (\prod_kS_{t_k})^{1/64}$ is lognormal and its call has a closed form (Kemna and Vorst, 1990): 6.1793. Its payoff has correlation 0.99941 with the
arithmetic payoff's, so $1 - \rho^2 = 0.0012$ and, with $\beta = 1.04$ estimated from the same paths, the variance falls by a factor of 853 at the same number of paths. Antithetic pairs do much less:
the two payoffs of a pair have correlation $-0.43$, a factor of 1.8 at equal payoff evaluations. Stratification is the idea behind the next section: it removes the variance between strata, and a
low-discrepancy point set stratifies every coordinate, and many of their products, at once.

\begin{definition}[Importance sampling]\label{def:qm:monte-carlo:is}
\emph{Importance sampling}\index{importance sampling} draws $X$ from a density $g$ instead of $f$ and averages $h(X)f(X)/g(X)$, the payoff times the likelihood ratio (a \omterm{def:qm:probability-at-speed:change}{Radon--Nikodym derivative});
it is unbiased whenever $g > 0$ where $hf \ne 0$.
\end{definition}

\begin{proposition}[Zero-variance importance density]\label{prop:qm:monte-carlo:zero}
For $h \ge 0$ with $\mu = \E_f[h(X)] > 0$, the density $g^* = hf/\mu$ makes $h(X)f(X)/g^*(X) = \mu$ for every draw: the \omterm{def:qm:estimation:estimator}{estimator} has zero variance.
\end{proposition}

\begin{proof}
Substitute: $hf/g^* = hf\mu/(hf) = \mu$ wherever $h > 0$, and $g^*$ puts no mass where $h = 0$.
\end{proof}

The optimal density needs the answer, but it says where to look: sample in proportion to where the payoff lives. For a digital paying if $S_T > 250$ (probability $1.21 \times 10^{-4}$, the
normal driver above $z^* = 3.67$), 100\,000 plain draws produce six hits and a relative \omterm{def:qm:estimation:estimator}{standard error} of 20\%. Shifting the driver's mean to $\theta$ and weighting each draw by $e^{-\theta Z +
\theta^2/2}$ brings it to 0.64\% at $\theta = 3.8$, a variance factor of about 1\,000 (\cref{fig:qm:monte-carlo:is}); shifting too far makes it worse again, as the weights of the few draws that
matter explode.

\begin{omfigure}
\begin{tikzpicture}
\begin{semilogyaxis}[width=10cm,height=4.8cm,xlabel={mean shift $\theta$ of the normal driver},ylabel={relative standard error},
  tick label style={font=\small},label style={font=\small},xmin=0,xmax=5,ymin=3e-3,ymax=0.5,grid=major,grid style={gray!25},table/col sep=comma]
\addplot[omDef,very thick,mark=*,mark size=1.4pt] table[x=theta,y=relse]{figdata/methods/26-monte-carlo/is.csv};
\draw[omThm,dashed,thick] (axis cs:3.67,3e-3) -- (axis cs:3.67,0.5) node[pos=0.85,right,font=\footnotesize]{$z^* = 3.67$};
\end{semilogyaxis}
\end{tikzpicture}

\omcaption{\omterm{def:qm:monte-carlo:is}{Importance sampling} for a digital paying if $S_T > 250$ ($S_0 = 100$, $\sigma = 25\%$, one year; probability $1.21 \times 10^{-4}$): relative \omterm{def:qm:estimation:estimator}{standard error} of 100\,000 draws against the
mean shift of the normal driver. Data: the chapter's tutorial, seeded.}\label{fig:qm:monte-carlo:is}
\end{omfigure}

\section{Low-discrepancy sequences}

\begin{definition}[Low-discrepancy sequence, Sobol sequence, quasi-Monte Carlo]\label{def:qm:monte-carlo:qmc}
The star discrepancy $D_n^*$ of points $x_1, \dots, x_n$ in $[0, 1)^d$ is the largest difference, over boxes $[0, a)$, between the fraction of points in the box and its volume. A \emph{low-discrepancy
sequence}\index{low-discrepancy sequence} has $D_n^* = O((\log n)^d/n)$. The \emph{Sobol sequence}\index{Sobol sequence} (Sobol, 1967) is a base-2 low-discrepancy sequence built coordinate by
coordinate from primitive polynomials and direction numbers, by exclusive-or of binary digits. \emph{Quasi-Monte Carlo}\index{quasi-Monte Carlo} replaces the random draws of Monte Carlo by such
points.
\end{definition}

\begin{proposition}[Koksma--Hlawka inequality]\label{prop:qm:monte-carlo:kh}
For $f$ of bounded variation $V(f)$ in the sense of Hardy and Krause, $\bigl|\frac1n\sum_if(x_i) - \int_{[0,1)^d}f\bigr| \le V(f)\,D_n^*$.
\end{proposition}

The bound gives \omterm{def:qm:monte-carlo:qmc}{quasi-Monte Carlo} a rate close to $1/n$ for smooth integrands of moderate effective dimension, but it is a worst-case bound that cannot be estimated from the points. Randomisation
restores an error estimate.

\begin{definition}[Randomised quasi-Monte Carlo, Brownian bridge construction]\label{def:qm:monte-carlo:rqmc}
\emph{Randomised quasi-Monte Carlo}\index{randomised quasi-Monte Carlo} randomises a low-discrepancy set so that each point is uniform while the set keeps its structure, and estimates the error
from independent replicates. The kit uses Owen's nested uniform scrambling (1995): each binary digit of each coordinate is flipped by a random bit attached to the node of the binary tree that its
leading digits define. The \emph{Brownian bridge construction}\index{Brownian bridge construction} assigns the first coordinate to $W_T$, the next to $W_{T/2}$ given its ends, and so on
down the levels, so that the first coordinates, where low-discrepancy points are best, carry most of the variance of the path.
\end{definition}

\Cref{fig:qm:monte-carlo:points} shows why: 256 scrambled Sobol points in two dimensions put exactly one point in each of 256 dyadic boxes of every shape, where independent draws leave gaps
and clusters. On the average-price call with $2^{14}$ paths and 32 independent scrambles, scrambled Sobol points with the increments in time order reduce the variance by a factor of about 34; with
the \omterm{def:qm:monte-carlo:rqmc}{Brownian bridge construction} by about 2\,500; with the bridge and the \omterm{def:qm:monte-carlo:vr}{control variate} together by about 37\,000, a \omterm{def:qm:estimation:estimator}{standard error} 190 times smaller at the same number of paths. The error also
falls faster than $n^{-1/2}$: fitted slopes between $-0.7$ and $-0.8$ from $2^6$ to $2^{15}$ paths (\cref{fig:qm:monte-carlo:cost}).

\begin{omfigure}
\begin{tikzpicture}
\begin{axis}[name=l,width=5.6cm,height=5.6cm,title={Philox},title style={font=\small},xmin=0,xmax=1,ymin=0,ymax=1,
  tick label style={font=\scriptsize},xtick={0,0.5,1},ytick={0,0.5,1},table/col sep=comma]
\addplot[only marks,mark=*,mark size=0.9pt,omThm] table[x=px,y=py]{figdata/methods/26-monte-carlo/points.csv};
\end{axis}
\begin{axis}[at={(l.east)},anchor=west,xshift=1cm,width=5.6cm,height=5.6cm,title={scrambled Sobol},title style={font=\small},xmin=0,xmax=1,ymin=0,ymax=1,
  tick label style={font=\scriptsize},xtick={0,0.5,1},ytick={0,0.5,1},table/col sep=comma]
\addplot[only marks,mark=*,mark size=0.9pt,omDef] table[x=sx,y=sy]{figdata/methods/26-monte-carlo/points.csv};
\end{axis}
\end{tikzpicture}

\omcaption{256 points in the unit square: independent Philox uniforms and the first two coordinates of Owen-scrambled Sobol points. Data: the chapter's tutorial,
seeded.}\label{fig:qm:monte-carlo:points}
\end{omfigure}

\begin{omfigure}
\begin{tikzpicture}
\begin{loglogaxis}[width=11cm,height=6cm,xlabel={number of paths},ylabel={standard error of the price},
  tick label style={font=\small},label style={font=\small},xmin=50,xmax=45000,ymin=1e-4,ymax=3,grid=major,grid style={gray!25},
  legend columns=3,legend cell align=left,legend style={font=\footnotesize,at={(0.5,-0.25)},anchor=north,draw=none,fill=none,
  /tikz/every even column/.append style={column sep=6pt}},table/col sep=comma]
\addplot[omThm,very thick,mark=*,mark size=1.4pt] table[x=n,y=plain]{figdata/methods/26-monte-carlo/cost.csv};
\addplot[omThm!50,very thick,mark=o,mark size=1.6pt] table[x=n,y=rqmcinc]{figdata/methods/26-monte-carlo/cost.csv};
\addplot[black,very thick,mark=square*,mark size=1.4pt] table[x=n,y=cv]{figdata/methods/26-monte-carlo/cost.csv};
\addplot[omDef!60,very thick,mark=triangle*,mark size=1.8pt] table[x=n,y=rqmcbridge]{figdata/methods/26-monte-carlo/cost.csv};
\addplot[omDef,very thick,mark=diamond*,mark size=1.8pt] table[x=n,y=rqmccv]{figdata/methods/26-monte-carlo/cost.csv};
\legend{plain Monte Carlo,{RQMC, time order},control variate,{RQMC, bridge},{RQMC, bridge + control}}
\end{loglogaxis}
\end{tikzpicture}

\omcaption{\omterm{def:qm:estimation:estimator}{Standard error} of the average-price call (64 fixings) against the number of paths: plain Monte Carlo and the \omterm{def:qm:monte-carlo:vr}{control variate} (slope $-\frac12$), and scrambled Sobol points (16
independent scrambles) with increments in time order, with the \omterm{def:qm:monte-carlo:rqmc}{Brownian bridge construction}, and with the bridge and the \omterm{def:qm:monte-carlo:vr}{control variate}. Data: the chapter's tutorial,
seeded.}\label{fig:qm:monte-carlo:cost}
\end{omfigure}

\section{Discretisation and multilevel methods}

\begin{definition}[Strong and weak orders of convergence, Milstein scheme]\label{def:qm:monte-carlo:order}
A scheme with step $h$ has \emph{strong order of convergence}\index{strong order of convergence} $\gamma$ if $\E|X_T^h - X_T| \le Ch^\gamma$, and \emph{weak order of convergence}\index{weak order of
convergence} $\beta$ if $|\E g(X_T^h) - \E g(X_T)| \le Ch^\beta$ for smooth $g$. The \emph{Milstein scheme}\index{Milstein scheme} for $dX = a(X)\,dt + b(X)\,dW$ adds to the Euler--Maruyama step the
term $\frac12b(X)b'(X)((\Delta W)^2 - h)$.
\end{definition}

\begin{proposition}[Orders of Euler and Milstein]\label{prop:qm:monte-carlo:orders}
Under Lipschitz and growth conditions (and smoothness of $a$, $b$ for Milstein), the \omterm{def:qm:stochastic-differential-equations:euler}{Euler--Maruyama scheme} has strong order $\frac12$ and weak order 1, and the \omterm{def:qm:monte-carlo:order}{Milstein scheme} strong order 1
(Kloeden and Platen, 1992).
\end{proposition}

\begin{proof}[Idea]
The Euler step misses the double stochastic integral $\int\!\!\int dW\,dW = \frac12((\Delta W)^2 - h)$, of size $h$ with mean zero, times $bb'$; these errors add like a random walk over $1/h$
steps, $h\sqrt{1/h} = h^{1/2}$ pathwise, while their means cancel, leaving an $O(h)$ bias in law. Milstein keeps that term, and the next missing terms are of order $h^{3/2}$ per step.
\end{proof}

A price needs the weak order; a hedge simulated along the path, or a multilevel \omterm{def:qm:estimation:estimator}{estimator}, needs the strong one. On a \omterm{def:qm:stochastic-differential-equations:gbm}{geometric Brownian motion} with $\sigma = 50\%$, driven by the same increments as the
exact solution, the fitted strong slopes are 0.49 for Euler and 0.89 for Milstein, and the weak error of the mean is exactly $X_0|(1 + \mu h)^n - e^{\mu T}|$, slope 1.00
(\cref{fig:qm:monte-carlo:orders}).

\begin{omfigure}
\begin{tikzpicture}
\begin{loglogaxis}[width=10cm,height=5cm,xlabel={time step $h$ (years)},ylabel={error},x dir=reverse,
  tick label style={font=\small},label style={font=\small},xmin=0.003,xmax=0.6,ymin=3e-4,ymax=20,grid=major,grid style={gray!25},
  legend columns=3,legend cell align=left,legend style={font=\footnotesize,at={(0.5,-0.3)},anchor=north,draw=none,fill=none,
  /tikz/every even column/.append style={column sep=8pt}},table/col sep=comma]
\addplot[omThm,very thick,mark=*,mark size=1.4pt] table[x=h,y=euler]{figdata/methods/26-monte-carlo/orders.csv};
\addplot[omDef,very thick,mark=square*,mark size=1.4pt] table[x=h,y=milstein]{figdata/methods/26-monte-carlo/orders.csv};
\addplot[black,very thick,mark=triangle*,mark size=1.8pt] table[x=h,y=weak]{figdata/methods/26-monte-carlo/orders.csv};
\legend{strong error (Euler),strong error (Milstein),weak error of the mean}
\end{loglogaxis}
\end{tikzpicture}

\omcaption{Strong error $\E|X_T^h - X_T|$ of the Euler and \omterm{def:qm:monte-carlo:order}{Milstein schemes} (20\,000 paths of a \omterm{def:qm:stochastic-differential-equations:gbm}{geometric Brownian motion}, $\mu = 5\%$, $\sigma = 50\%$, $X_0 = 100$, one year) and the exact weak error of
the mean, against the time step. Data: the chapter's tutorial, seeded.}\label{fig:qm:monte-carlo:orders}
\end{omfigure}

\begin{definition}[Multilevel Monte Carlo]\label{def:qm:monte-carlo:mlmc}
\emph{Multilevel Monte Carlo}\index{multilevel Monte Carlo} (Giles, 2008) writes the finest-level expectation as a telescoping sum $\E P_L = \E P_0 + \sum_{l=1}^L\E[P_l - P_{l-1}]$, where $P_l$ is the
payoff computed with step $h_l = 2^{-l}T$, and estimates each term independently, with $P_l$ and $P_{l-1}$ in each correction computed from the same Brownian path.
\end{definition}

\begin{proposition}[Complexity of multilevel Monte Carlo]\label{prop:qm:monte-carlo:giles}
If the bias of $P_l$ is $O(2^{-\alpha l})$, the variance of the correction $V_l = O(2^{-\beta l})$ and its cost $C_l = O(2^{\gamma l})$, with $\alpha \ge \frac12\min(\beta, \gamma)$, then a mean square error
$\epsilon^2$ costs $O(\epsilon^{-2})$ if $\beta > \gamma$, $O(\epsilon^{-2}(\log\epsilon)^2)$ if $\beta = \gamma$ and $O(\epsilon^{-2 - (\gamma - \beta)/\alpha})$ if $\beta < \gamma$; standard Monte Carlo costs
$O(\epsilon^{-2 - \gamma/\alpha})$.
\end{proposition}

\begin{proof}[Idea]
Minimising the cost $\sum_lN_lC_l$ subject to the variance constraint $\sum_lV_l/N_l = \epsilon^2/2$ gives $N_l \propto\sqrt{V_l/C_l}$ and a total cost
\[2\epsilon^{-2}\Bigl(\sum_l\sqrt{V_lC_l}\Bigr)^2;\]
the sum is dominated by the coarsest levels
if $\beta > \gamma$, by the finest if $\beta < \gamma$, and grows like $L \sim \log\epsilon^{-1}$ if they are equal. The finest level $L$ is set by the bias, $2^{-\alpha L} \sim \epsilon$.
\end{proof}

The strong order sets $\beta$: Milstein on a Lipschitz payoff gives $\beta = 2 > \gamma = 1$, and the corrections' variances in the tutorial fall by a factor of about four per level
(0.23, 0.060, 0.016, 0.0045, 0.0012, 0.00030, 0.000079). For a European call ($\alpha = 1$), the driver listed in \cref{tut:qm:monte-carlo:asian} reaches a root mean square error of 0.005 with 7 levels
at a cost of $3.2 \times 10^7$ time steps, where standard Monte Carlo at the same finest step would need $3.0 \times 10^9$: a factor of 94, which grows as $\epsilon$ falls (8, 20, 45, 94 for
$\epsilon$ from 0.04 to 0.005; \cref{fig:qm:monte-carlo:mlmc}).

\begin{omfigure}
\begin{tikzpicture}
\begin{loglogaxis}[width=10cm,height=5cm,xlabel={target root mean square error $\epsilon$},ylabel={cost (time steps)},x dir=reverse,
  tick label style={font=\small},label style={font=\small},xmin=0.004,xmax=0.05,ymin=1e5,ymax=1e10,grid=major,grid style={gray!25},
  xtick={0.005,0.01,0.02,0.04},xticklabels={0.005,0.01,0.02,0.04},
  legend columns=2,legend cell align=left,legend style={font=\footnotesize,at={(0.5,-0.3)},anchor=north,draw=none,fill=none,
  /tikz/every even column/.append style={column sep=8pt}},table/col sep=comma]
\addplot[omThm,very thick,mark=*,mark size=1.6pt] table[x=eps,y=standard]{figdata/methods/26-monte-carlo/mlmc.csv};
\addplot[omDef,very thick,mark=square*,mark size=1.6pt] table[x=eps,y=mlmc]{figdata/methods/26-monte-carlo/mlmc.csv};
\legend{standard Monte Carlo (Milstein),multilevel Monte Carlo}
\end{loglogaxis}
\end{tikzpicture}

\omcaption{Cost of pricing a European call ($S_0 = K = 100$, $r = 3\%$, $\sigma = 25\%$, one year) to a target root mean square error: \omterm{def:qm:monte-carlo:mlmc}{multilevel Monte Carlo} with Milstein corrections against standard
Monte Carlo at the same finest step. Data: the chapter's tutorial, seeded.}\label{fig:qm:monte-carlo:mlmc}
\end{omfigure}

\section{Tutorial: sixteen times fewer machines}\label{tut:qm:monte-carlo:asian}

\begin{tutorial}
\textbf{Goal.} Value the batch's representative trade five ways and measure each \omterm{def:qm:estimation:estimator}{estimator}'s variance at equal paths. \textbf{End state:} the factors 1.8, 853, 34, 2\,500 and 37\,000 and
\cref{fig:qm:monte-carlo:cost}.
\begin{tutsteps}
\item \textbf{Counter-based uniforms.} Philox4x64-10 in C++20, reproducing NumPy's \texttt{Philox} block for block.
\omcode{code/firm/mcengine/cpp/firm_mcengine.hpp}{113}{124}{Philox4x64-10: ten rounds of two wide multiplications (C++20).}\label{lst:qm:monte-carlo:philox}
\item \textbf{Scrambled Sobol points.} Gray-code Sobol points from Joe and Kuo's direction numbers, and Owen's scrambling by hashed tree nodes.
\omcode{code/firm/mcengine/firm_mcengine.py}{268}{297}{Sobol points and nested uniform scrambling (Python).}\label{lst:qm:monte-carlo:sobol}
\item \textbf{\omterm{def:qm:estimation:estimator}{Estimators}.} Run \texttt{table()} in \texttt{qm\_mc.py}: plain, antithetic, \omterm{def:qm:monte-carlo:vr}{control variate}, and scrambled Sobol with and without the bridge (32 scrambles of $2^{14}$ points); then
\texttt{error\_vs\_cost()}.
\item \textbf{Multilevel.} Run \texttt{mlmc\_study()}; the driver is below.
\omcode{code/firm/mcengine/firm_mcengine.py}{384}{412}{Giles's multilevel Monte Carlo driver (Python).}\label{lst:qm:monte-carlo:mlmc}
\end{tutsteps}
\textbf{What to change next.} Put the strike at 130 and watch the \omterm{def:qm:monte-carlo:vr}{control variate}'s factor fall with the correlation; replace the bridge by a principal-component construction; price a barrier with
the bridge crossing probability of chapter~2 inside the multilevel driver.
\end{tutorial}

\section{Build: the Monte Carlo engine, stage three}\label{bld:qm:monte-carlo:mcengine}

\begin{build}
\bfield{Purpose} Make each path worth more and every path reproducible: the generators, point sets and \omterm{def:qm:estimation:estimator}{estimators} that the firm's pricing and risk batches call.
\bfield{Interface} \texttt{philox4x64}, \texttt{philox\_uniforms(n, dim, seed, stream, first\_path)}; \texttt{sobol\_points}, \texttt{owen\_scramble}; \texttt{norm\_ppf};
\texttt{bridge\_from\_normals}; \texttt{asian\_payoffs}, \texttt{geometric\_asian\_price}, \texttt{control\_variate}; \texttt{mlmc(sampler, eps)}. C++20 and Rust twins of \texttt{philox4x64},
\texttt{sobol\_point} (eight dimensions) and \texttt{owen\_scramble}.
\bfield{Rules} Counters, not sequential state, identify paths; every estimate is reported with its \omterm{def:qm:estimation:estimator}{standard error} (from replicates for \omterm{def:qm:monte-carlo:rqmc}{randomised quasi-Monte Carlo}, never from within one point
set); unscrambled Sobol points are never used for error bars; direction numbers carry their licence.
\bfield{Acceptance tests} \texttt{code/firm/mcengine/\{tests,cpp,rust\}}: Philox against NumPy's; splitting paths reproduces them; Joe and Kuo's published first points; identical scrambled integers in
the three languages; the net property after scrambling; $\Phi^{-1}$ against the standard library to $10^{-8}$; bridge covariance; the geometric closed form against simulation and against
Black--Scholes for one fixing; the multilevel driver on a sampler with a known answer.
\bfield{Stretch} Joe and Kuo's full 21\,201 dimensions; a principal-component path construction; conditional Monte Carlo for barriers; adjoint Greeks along the paths (chapter~28).
\end{build}

\begin{omsources}
\item N.~Metropolis and S.~Ulam, ``The Monte Carlo method'', \emph{Journal of the American Statistical Association} 44, 1949; P.~P.~Boyle, ``Options: a Monte Carlo approach'', \emph{Journal of
Financial Economics} 4, 1977.
\item P.~Glasserman, \emph{Monte Carlo Methods in Financial Engineering}, Springer, 2003.
\item J.~K.~Salmon, M.~A.~Moraes, R.~O.~Dror and D.~E.~Shaw, ``Parallel random numbers: as easy as 1, 2, 3'', SC11, 2011; M.~Matsumoto and T.~Nishimura, \emph{ACM TOMACS} 8, 1998; M.~E.~O'Neill,
PCG, Harvey Mudd College report HMC-CS-2014-0905, 2014.
\item I.~M.~Sobol', \emph{USSR Computational Mathematics and Mathematical Physics} 7, 1967; S.~Joe and F.~Y.~Kuo, \emph{SIAM Journal on Scientific Computing} 30, 2008 (direction numbers
\texttt{new-joe-kuo-6.21201}); A.~B.~Owen, ``Randomly permuted $(t, m, s)$-nets and $(t, s)$-sequences'', 1995; B.~Moskowitz and R.~E.~Caflisch, \emph{Mathematical and Computer Modelling} 23, 1996.
\item A.~G.~Z.~Kemna and A.~C.~F.~Vorst, \emph{Journal of Banking and Finance} 14, 1990.
\item M.~B.~Giles, ``Multilevel Monte Carlo path simulation'', \emph{Operations Research} 56, 2008; P.~E.~Kloeden and E.~Platen, \emph{Numerical Solution of Stochastic Differential Equations},
Springer, 1992; G.~N.~Milstein, \emph{Theory of Probability and its Applications} 19, 1975.
\item P.~J.~Acklam, ``An algorithm for computing the inverse normal cumulative distribution function'' (web note, 2004).
\end{omsources}

\section{Exercises}

\begin{exercise}[$\star$]\label{exo:qm:monte-carlo:1}
A batch has \omterm{def:qm:estimation:estimator}{standard error} 80\,000 dollars with 100\,000 paths. How many paths reach 20\,000, and what variance-reduction factor does the same at 100\,000 paths?
\end{exercise}

\begin{exercise}[$\star$]\label{exo:qm:monte-carlo:2}
A \omterm{def:qm:monte-carlo:vr}{control variate} has correlation 0.9 with the payoff. What is the variance factor? And with 0.99?
\end{exercise}

\begin{exercise}[$\star$]\label{exo:qm:monte-carlo:3}
Why do \omterm{def:qm:monte-carlo:vr}{antithetic variates} help a monotone payoff and not an even function of $Z$ such as $Z^2$?
\end{exercise}

\begin{exercise}[$\star\star$]\label{exo:qm:monte-carlo:4}
Show that proportional stratification never increases the variance: with strata $s$ of probability $p_s$ and conditional variances $\sigma_s^2$, the stratified variance is $\frac1n\sum_sp_s\sigma_s^2 \le
\sigma^2/n$.
\end{exercise}

\begin{exercise}[$\star\star$]\label{exo:qm:monte-carlo:5}
For $P(Z > z^*)$ with $Z$ standard normal, compute the second moment of the importance-sampling \omterm{def:qm:estimation:estimator}{estimator} with the shift $\theta$, and show that it is minimised near $\theta = z^*$ for large $z^*$.
\end{exercise}

\begin{exercise}[$\star\star$]\label{exo:qm:monte-carlo:6}
Derive the Milstein step for \omterm{def:qm:stochastic-differential-equations:gbm}{geometric Brownian motion} and show that it is the second-order expansion of the exact step $X e^{(\mu - \sigma^2/2)h + \sigma\Delta W}$ in $\Delta W$.
\end{exercise}

\begin{exercise}[$\star\star\star$]\label{exo:qm:monte-carlo:7}
\emph{Coding.} Compute the variance factor of the \omterm{def:qm:monte-carlo:vr}{control variate} at strikes 80, 100 and 130, and explain the trend.
\end{exercise}

\begin{exercise}[$\star\star\star$]\label{exo:qm:monte-carlo:8}
\emph{Find the flaw.} ``We moved the batch to Sobol points, 100\,000 per trade, and report the \omterm{def:qm:estimation:estimator}{standard error} $s/\sqrt n$ from the points as before; it fell by a factor of ten.''
\end{exercise}

\section{Problem: Sixteen Times Fewer Machines}

\begin{problem}[{Weekend problem --- a factor of four in error, without sixteen times the paths}]\label{pb:qm:monte-carlo:1}
The representative trade is an arithmetic-average call on 64 fixings ($S_0 = K = 100$, $r = 3\%$, $\sigma = 25\%$, one year), valued with $2^{14}$ paths; the batch needs a factor of four in \omterm{def:qm:estimation:estimator}{standard error}.

\medskip\noindent\textbf{Part I --- The plain \omterm{def:qm:estimation:estimator}{estimator}.}
\begin{enumerate}
\item What are the plain estimate, its \omterm{def:qm:estimation:estimator}{standard error} and the reference price?
\item What variance factor does the target require?
\item What does the cost-variance product say about comparing \omterm{def:qm:estimation:estimator}{estimators}?
\item Why does the \omterm{def:qm:monte-carlo:prng}{counter-based generator} make the batch reproducible across machines?
\item Why is the inverse transform, rather than Box--Muller, used to feed low-discrepancy points?
\end{enumerate}

\medskip\noindent\textbf{Part II --- \omterm{def:qm:monte-carlo:mc}{Variance reduction}.}
\begin{enumerate}[resume]
\item What factor do antithetic pairs give, and why so little?
\item What is the \omterm{def:qm:monte-carlo:vr}{control variate}, its closed-form value, its correlation and its factor?
\item What is the optimal coefficient, and what was estimated?
\item Where does \omterm{def:qm:monte-carlo:is}{importance sampling} help, and by how much on the deep out-of-the-money digital?
\item Why does too large a shift make it worse?
\end{enumerate}

\medskip\noindent\textbf{Part III --- Points and paths.}
\begin{enumerate}[resume]
\item What factor do scrambled Sobol points give with the increments in time order, and with the bridge?
\item Why does the bridge matter so much?
\item What do the two together give, and what is the error's slope in the number of paths?
\item How is the error of \omterm{def:qm:monte-carlo:rqmc}{randomised quasi-Monte Carlo} estimated?
\item What do the strong and weak orders measure, and what are their fitted values here?
\end{enumerate}

\medskip\noindent\textbf{Part IV --- Judgement.}
\begin{enumerate}[resume]
\item What does \omterm{def:qm:monte-carlo:mlmc}{multilevel Monte Carlo} save on the European call at $\epsilon = 0.005$, and why?
\item Would the factor of 37\,000 hold across the batch?
\item What must the batch log so that one trade's paths can be regenerated?
\item State the \emph{named result}: the variance-reduction factor of the \omterm{def:qm:monte-carlo:vr}{control variate} plus \omterm{def:qm:monte-carlo:rqmc}{randomised quasi-Monte Carlo} on the representative trade, against the 16 required.
\item In one sentence: where does the work of a Monte Carlo engine go?
\end{enumerate}
\end{problem}

\section{Interview questions}

\begin{interviewq}[$\star$ \iqroles{researcher, developer}]\label{iq:qm:monte-carlo:1}
How does the Monte Carlo error scale with the number of paths and with the dimension?
\end{interviewq}

\begin{interviewq}[$\star\star$ \iqroles{researcher}]\label{iq:qm:monte-carlo:2}
Name three variance-reduction techniques and when each works.
\end{interviewq}

\begin{interviewq}[$\star\star$ \iqroles{developer}]\label{iq:qm:monte-carlo:3}
How do you make a distributed Monte Carlo batch reproducible regardless of the number of machines?
\end{interviewq}

\begin{interviewq}[$\star\star$ \iqroles{researcher}]\label{iq:qm:monte-carlo:4}
Why might \omterm{def:qm:monte-carlo:qmc}{quasi-Monte Carlo} do badly on a 250-step path, and what fixes it?
\end{interviewq}

\begin{interviewq}[$\star\star$ \iqroles{researcher, risk}]\label{iq:qm:monte-carlo:5}
How would you estimate the probability of a loss beyond a very high threshold by simulation?
\end{interviewq}

\begin{interviewq}[$\star\star\star$ \iqroles{researcher}]\label{iq:qm:monte-carlo:6}
Explain \omterm{def:qm:monte-carlo:mlmc}{multilevel Monte Carlo} and why the strong order of the scheme matters for it.
\end{interviewq}
